\section{Introducción}

Cadenas de suministro inteligentes integran IoT, blockchain y análisis de datos para monitoreo en tiempo real y optimización. En el agro colombiano, esto significa menos pérdidas y mayor eficiencia.

El Portal Agro Comercial del Huila conecta productores, transportistas y compradores, creando cadenas inteligentes que responden dinámicamente a necesidades.

\section{Marco teórico y trabajos relacionados}

Estudios sobre supply chains inteligentes (2024) destacan beneficios en reducción de costos. Investigaciones en blockchain para trazabilidad (2024) y IoT en logística (2023).

Trabajos de la FAO (2023) sobre cadenas de valor agroalimentarias.

\section{Metodología}

\subsection{Diseño de Cadenas}
Análisis de flujos actuales y diseño de sistemas inteligentes.

\subsection{Implementación Tecnológica}
Integración de sensores y plataformas de seguimiento.

\subsection{Evaluación}
Medición de eficiencia y reducción de pérdidas.

\section{Implementación del Portal Agro Comercial del Huila}

Módulos para coordinación logística y trazabilidad.

\subsection{Conexión de Actores}
Plataforma para contratos y seguimiento.

\subsection{Herramientas Inteligentes}
Análisis predictivo para optimización.

\section{Resultados e Impactos}

Reducción de pérdidas en un 25\% y mejora en tiempos de entrega.

\subsection{Eficiencia Logística}
Mejor coordinación reduce costos.

\subsection{Sostenibilidad}
Menos desperdicio y huella ambiental.

\section{Discusión}

\subsection{Desafíos}
Integración de actores diversos.

\subsection{Comparación con Literatura}
Resultados positivos, alineados con tendencias globales.

\section{Conclusiones y Visión Futura}

Cadenas inteligentes son el futuro del agro. El portal facilita su adopción.

\subsection{Trabajo Futuro}
Integrar más tecnologías emergentes.

\section{Implementación en el proyecto Portal Agrocomercial del Huila}

El Portal Agrocomercial del Huila crea cadenas de suministro inteligentes mediante módulos de trazabilidad con blockchain, seguimiento en tiempo real con IoT y análisis predictivo para optimización logística. Permite a productores coordinar transporte, gestionar inventarios y responder a demanda dinámica, reduciendo pérdidas y costos. Además, el portal fomenta colaboraciones entre actores, fortaleciendo la resiliencia y sostenibilidad de las cadenas agroalimentarias en el Huila.